\section{Coupled population and particle PWGF}
\label{sec:pwgf-coupling}

We define the population and finite-particle dynamics and the coupling used to
compare them along a prescribed population PWGF path.

\subsection{Population WGF and Hessian-guided perturbations}
\label{subsec:population-pwgf}

Let \(F:\operatorname{dom}F\to\R\) be the objective functional.
Between perturbation times, the
population curve \((\mu_t)_{t\ge 0}\) follows the Wasserstein gradient flow
\begin{equation}
    \partial_t \mu_t + \nabla\cdot(v_t\mu_t)=0,
    \qquad
    v_t(x):=-\dW F(\mu_t,x).
    \label{eq:population-wgf}
\end{equation}
The original population PWGF triggers perturbations adaptively when
\(\|\dW F(\mu_t)\|_{L^2(\mu_t)}\) falls below a prescribed first-order threshold;
our finite-sample comparison uses deterministic jump times
$0<t_1<\cdots<t_{\NofJ}\le T$, fixed in advance and common to both systems.
At each such time, \citet{yamamoto2025hessianguided} perturb the measure by pushing it forward along a Hessian-guided Gaussian vector field.  The idea parallels perturbed
gradient methods in Euclidean space~\citep{GeHuangJinYuan2015,pmlr-v75-jin18a}: near a saddle,
randomness should reach negative-curvature directions.  In
\(\mathcal P_2(\mathbb R^d)\), (i) the definition of such perturbations and (ii) their
alignment with descent directions are both nontrivial, motivating the
Hessian-induced covariance.
In Wasserstein space, the second variation
along the pushforward \((\Id+hv)_\#\mu\) is represented by
\(\mathcal H_\mu=H_\mu+H_\mu'\), where
\[
    (H_\mu v)(x):=\int \HessW F(\mu;x,y)v(y)\,\mu(dy),
    \qquad
    (H_\mu'v)(x):=\xHessW F(\mu,x)v(x).
\]
See \citet{ambrosio2005gradient} for detailed properties. The perturbation studied here uses the squared integral
Hessian $H_\mu^2$ as covariance. The tracking curvature uses the full operator
$\mathcal H_\mu$; first-order stationarity alone does not imply $H_\mu'=0$.
The covariance kernel of \(H_\mu^2\) is
\[
    K_\mu(x,y)
    :=
    \int
        \HessW F(\mu;x,z)\HessW F(\mu;z,y)\,\mu(\dd z).
    \label{eq:hessian-kernel}
\]
This covariance preserves the eigenspaces of $H_\mu$ and assigns variance proportional to its squared eigenvalues. Relating this covariance to descent directions of the full Hessian requires additional structural assumptions. The population jump is
\begin{equation}
    \mu_{t_\ell}
    =
    (\Id+\eta_\ell \xi_{\mu_{t_\ell-}})_\#\mu_{t_\ell-},
    \qquad
    \xi_{\mu_{t_\ell-}}\sim \GP(0,K_{\mu_{t_\ell-}}),
    \label{eq:population-jump}
\end{equation}
with perturbation size \(\eta_\ell>0\).

\subsection{Empirical particle dynamics}
\label{subsec:particle-pwgf}

For \(\bs{X}=(X^1,\ldots,X^N)\in(\R^d)^N\), write $\rho_{\bs{X}}\coloneq \frac1N\sum_{i=1}^N \delta_{X^i}.$
The empirical PWGF evolves by the particle ODE
\begin{equation}
    \frac{\dd}{\dd t}X^i_t
    =
    -\dW F(\rho_{\mathbf X_t},X^i_t),
    \qquad i=1,\ldots,N,
    \label{eq:particle-flow}
\end{equation}
between jump times.  At \(t_\ell\), the particles are updated by
\begin{equation}
    X^i_{t_\ell}
    =
    X^i_{t_\ell-}+\eta_\ell\xi_{\rho_{\mathbf X_{t_\ell-}}}(X^i_{t_\ell-}),
    \qquad i=1,\ldots,N,
    \label{eq:particle-jump}
\end{equation}
where \(\xi_{\rho_{\mathbf X_{t_\ell-}}} \sim \GP(0,K_{\rho_{\mathbf X_{t_\ell-}}})\) is a Gaussian process induced by the empirical Hessian-guided kernel.  The covariance of \(\xi_{\rho_{\mathbf X}}(X^1),\ldots,\xi_{\rho_{\mathbf X}}(X^N)\) is
a block matrix whose \((i,j)\)-block is \(K_{\rho_{\mathbf X}}(X^i,X^j)\).
For a fixed configuration, this jump is finite dimensional: one applies an
\(Nd\times Nd\) covariance factor to a standard Gaussian vector in
\(\mathbb R^{Nd}\). It can therefore be simulated without access to the
population path.

\subsection{The auxiliary reference particles}
\label{subsec:auxiliary-particles}

The comparison is cleaner if we introduce a reference particle system
\(Y_t=(Y^1_t,\ldots,Y^N_t)\). Conditional on the population fields,
\(Y_t^1,\ldots,Y_t^N\) are i.i.d. with law \(\mu_t\). The reference particles
start from the same initial positions as the empirical particles and evolve
with the population drift and jump covariance. Their Gaussian increments are
coupled to the empirical increments through the common Gaussian vector
described below.
Specifically, we set
\begin{equation}
    X^i_0=Y^i_0,
    \qquad
    Y^i_0\stackrel{\mathrm{i.i.d.}}\sim \mu_0.
    \label{eq:same-initial-particles}
\end{equation}
Between jumps,
\begin{equation}
    \frac{\dd}{\dd t}Y^i_t
    =
    -\dW F(\mu_t,Y^i_t),
    \qquad i=1,\ldots,N.
    \label{eq:reference-flow}
\end{equation}
At \(t_\ell\), the reference jump is
\begin{equation}
    Y^i_{t_\ell}
    =
    Y^i_{t_\ell-}+\eta_\ell\xi_{\mu_{t_\ell-}}(Y^i_{t_\ell-}),
    \qquad i=1,\ldots,N,
    \label{eq:reference-jump}
\end{equation}

At each jump, the population Gaussian field is generated from a fresh
innovation independent of the initial sample and the previous innovations.
Conditional on these fields, the flow and jump maps are deterministic;
applying their common composition to the independent initial positions
preserves the reference particles' conditional i.i.d. law
(\Cref{lem:population-first-finite-coupling,eq:population-pushforward-product-law}).
For the comparison, \Cref{lem:population-first-finite-coupling} constructs a common
standard Gaussian vector \(\mathbf g_\ell^N\in\mathbb R^{Nd}\) from the given
population-field evaluations and independent auxiliary Gaussian noise,
leaving the population path unchanged. If
\(B_{\rm pop}\) denotes the aligned factor reproducing the population-field
values at the reference particles and \(M_X\) denotes the empirical Hessian
factor, then
\[
    \bigl(\xi_{\mu_{t_\ell-}}(Y^i_{t_\ell-})\bigr)_{i=1}^N
    =B_{\rm pop}\mathbf g_\ell^N,
    \qquad
    \bigl(\xi_{\rho_{\mathbf X_{t_\ell-}}}(X^i_{t_\ell-})\bigr)_{i=1}^N
    =M_X\mathbf g_\ell^N.
\]
Both jumps have the prescribed conditional laws, and the Gaussian
field-value vectors differ by \((M_X-B_{\rm pop})\mathbf g_\ell^N\).
The precise construction, including
singular covariances, and the conditional i.i.d. law of the reference particles
are given in
\Cref{lem:population-first-finite-coupling,lemma-gp-jump-implies-conditional-iid}.

\subsection{Tracking metric and decomposition}
\label{subsec:tracking-metric}

We measure the pathwise discrepancy by
$r_t
    \coloneq 
    \left(\frac1N\sum_{i=1}^N \|X^i_t-Y^i_t\|^2\right)^{1/2}
    =
    \frac1{\sqrt N}\|X_t-Y_t\|_2.
    \label{eq:tracking-metric}
$
Since the empirical coupling \(N^{-1}\sum_i\delta_{(X^i_t,Y^i_t)}\) is admissible, $\Wtwo(\rho_{\mathbf X_t},\rho_{\mathbf Y_t})\le r_t$.
To compare objective values with the population flow, we combine the
paired-particle estimate with concentration of the reference cloud's feature
moments (\Cref{subsec:functional-gap}).
The next section gives joint tracking and objective bounds under a prescribed
common schedule.
